% Generated draft: requires mathematical and source review.
\documentclass[11pt]{article}
\usepackage[margin=25mm]{geometry}
\usepackage{fontspec}
\setmainfont{Latin Modern Roman}
\usepackage{amsmath,amssymb,mathtools,unicode-math}
\setmathfont{Latin Modern Math}
\usepackage{enumitem,needspace,xurl,hyperref}
\setlength{\parindent}{0pt}
\setlength{\parskip}{5pt}
\setlength{\emergencystretch}{3em}
\newcommand{\N}{\mathbb{N}}
\newcommand{\Z}{\mathbb{Z}}
\newcommand{\Q}{\mathbb{Q}}
\newcommand{\R}{\mathbb{R}}
\newcommand{\C}{\mathbb{C}}
\newcommand{\F}{\mathbb{F}}
\newcommand{\A}{\mathbb{A}}
\newcommand{\PP}{\mathbb{P}}
\newcommand{\E}{\mathbb{E}}
\newcommand{\eps}{\varepsilon}
\newcommand{\dd}{\,\mathrm{d}}
\begin{document}
\section*{UnsolvedMath ID: 800175}
\textbf{Source problem code:} catalogue entry
\subsection{Definitions and mathematical statement}
Let \(\zeta(s)\) denote the meromorphic continuation of the Riemann zeta function, initially defined by \(\zeta(s)=\sum_{n=1}^{\infty}n^{-s}\) for \(\operatorname{Re}s>1\). For \(p>0\) and \(T\geq 2\), define the critical-line moment
\[
M_p(T)=\int_T^{2T}\left|\zeta\!\left(\tfrac12+it\right)\right|^p\,\dd t.
\]
To express the zero-density question, use the convention
\[
N(\sigma,T)=\#\left\{\rho=\beta+i\gamma:\zeta(\rho)=0,\ 0<\beta<1,\ \beta\geq\sigma,\ 0<\gamma\leq T\right\},
\qquad \tfrac12<\sigma<1,
\]
where zeros are counted with multiplicity. These counting endpoints and parameter ranges are conventions for this draft; the source does not specify them.

\textbf{Research problem (Timothy Trudgian).} Investigate quantitatively how upper bounds for \(M_p(T)\) can be used to obtain upper bounds for \(N(\sigma,T)\) as \(T\to\infty\), with the dependence on the moment estimate and on \(\sigma\) made explicit. The source states that Ingham's zero-density estimate uses \(M_4(T)\). Determine whether and how other nontrivial moment estimates, specifically including Heath-Brown's estimate for \(M_{12}(T)\), can be used in such zero-density results. Investigate analogous connections between moments and zero density for other \(L\)-functions.

\textbf{Unresolved scope.} The source poses an investigation rather than a single quantified conjecture. It supplies neither a target density bound nor the precise moment estimates to be used. It also does not specify a class of other \(L\)-functions, their moment integrals, or whether the intended density estimates concern individual functions or families. These choices must be stated in any precise theorem developed from the question; no universal implication from an arbitrary moment bound is asserted here.
\subsection{Short English statement}
Investigate how estimates for moments of the Riemann zeta function lead to estimates counting zeros near the line \(\operatorname{Re}s=1\). In particular, examine whether moments other than the fourth, including Heath-Brown's twelfth-moment estimate, can be used, and whether the method extends to other \(L\)-functions. The source leaves the desired bounds and the class of \(L\)-functions unspecified.
\subsection{Sources}
\begin{itemize}
\item Local UnsolvedMath registry, permanent ID \(800175\), source record \texttt{catalogue entry}, titled \emph{Moments of L-functions --- Moments of the Riemann zeta function --- Problem 26}. The supplied registry labels the problem open; this is a source label, not an independently verified status assessment.
\item Hamieh, Kadiri, and Ng, \emph{Moments of L-functions Problem List} (2024), Problem 26: \url{https://arxiv.org/abs/2405.17800}. The supplied provenance records extraction by \texttt{arxiv\_latex} and access on 2026-10-16. That access date is supplied provenance; no independent source visit or verification is claimed.
\item The problem text attributes the question to Timothy Trudgian, UNSW Canberra. The registry's \texttt{proposed\_by} field lists Hamieh, Kadiri, and Ng, and its proposed year is 2024; both attribution layers are retained.
\item The supplied statement names Ingham's zero-density estimate and Heath-Brown's twelfth-moment estimate, but supplies no individual bibliographic citations or exact formulations for these results.
\end{itemize}
\end{document}
